\scanpage{70}
5) Soient $\rho,\rho',\rho''$ les trois autres triades de $X$, qui sont donc mutuellement parallèles, alors $\overline\rho,\overline\rho',\overline\rho''$ sont \unsure{formés|formées} respectivement des éléments dont l'ombre sur $X$ est $\rho,\rho',\rho''$, et leur réunion est un pavé $X_2$, dont elles forment une famille de triades parallèles.
Alors $X_1,X_2$ sont les seuls pavés de $S$ contenus dans $S-X$, \deleted{et} \inserted{ils forment une \uline{partition} de $S\setminus X$,} l'ens \inserted{(de card 2)} de ces pavés est en corr. biun. can. avec l'ens \inserted{$\varepsilon$} des deux familles de triades de $X$.

6) L'\deleted{famille} \inserted{ens $I'_1$} des 3 triades de $X_1$ distinctes des $\overline\tau,\overline\tau',\overline\tau''$, et l'ens \inserted{$I'_2$} des 3 triades de $X_2$ distinctes de $\overline\rho,\overline\rho',\overline\rho''$, est en corr. $1$-$1$, \deleted{par} la relation d'opposition. \inserted{Les 6 triades forment un hexagone de triades liées, comme tantôt.}
Soit $I_0$ l'ens quotient (= ens des trois \inserted{diagonales de cet hexagone})

Pour tte triades \inserted{$\sigma$, élément} d'une des familles $\{\overline\tau,\overline\tau',\overline\tau''\}$ ou $(\overline\rho,\overline\rho',\overline\rho'')$, \inserted{i.e. \unsure{fibre} ; une fibre de $S-X\longrightarrow\mathrm{Tria}(X)$,} on a une bijection canonique $\sigma\simeq I_0$, \deleted{\unsure{résultant}} \inserted{induite} par les applic\textsuperscript{tions} can. $X_1\longrightarrow I_0$ et $X_2\longrightarrow I_0$ (qui donnent lieu : $X_1\overset{\sim}{\longrightarrow}I_1\times I_0,\ X_2=I_2\times I_0$)

\deleted{Les \unsure{où} $I_0$ est aussi l'ens des \unsure{$2$}$\times$triades contenues dans}

Pour que une triade $\sigma$ soit $\in I'_1\cup I'_2$, il f. et s. que toutes les triades de l'hexagone des triades liées la contenant, ou seulement la

% archive sheet number: 69

